\documentclass[11pt]{article}
\usepackage[margin=1in]{geometry}
\usepackage{hyperref}
\usepackage{booktabs}
\usepackage{enumitem}
\usepackage{graphicx}
\usepackage{float}

\title{Final Report: Agentic Workflow for Automated Binary/APK Analysis}
\author{Abdeljalil Otman}
\date{March 22, 2026}

\begin{document}
\maketitle

\begin{abstract}
This report describes a practical, agentic malware analysis workflow that automates analysis of binary executables and Android APKs. The system uses high-level, specialized tools exposed through Model Context Protocol (MCP) servers and orchestrated by an LLM-based agent (Agno). The design focuses on clear abstractions (not raw command wrappers), reproducible containerization, and step-by-step validation. I evaluate the workflow on small synthetic samples to validate end-to-end analysis under low-resource constraints.
\end{abstract}

\section{Introduction}
Malware analysis traditionally involves a sequence of specialized tools and significant manual effort to interpret results. Modern malware employs sophisticated obfuscation techniques such as runtime code unpacking and polymorphic encoding to evade detection. Binary similarity detection, as explored in recent neural network approaches, offers scalable alternatives to manual signature generation. Understanding these techniques is critical for building automated analysis systems that can adapt to evolving malware families.

As a student, I wanted a workflow that still feels rigorous but reduces repetitive steps while maintaining transparency. This project builds an agentic pipeline that coordinates high-level analysis tools through an LLM-based agent (Agno). The motivation stems from: (1) the complexity of modern malware obfuscation requiring multiple specialized tools, (2) the need for clear tool abstractions that work well with agent reasoning, (3) MCP-based interoperability to reduce tool coupling, and (4) low-cost execution via the OpenRouter free tier to support resource-constrained environments.

The design draws inspiration from neural approaches to binary analysis~\cite{gemini2018} while implementing practical extraction of hidden code in packed executables~\cite{polyunpack2007}. Rather than attempting full graph embedding or heavy instrumentation, this implementation focuses on lightweight heuristics that remain effective on low-end hardware while exposing meaningful signals to an LLM coordinator.

\section{System Overview}
\subsection{Architecture}
The system architecture is designed around separable analysis stages, allowing the LLM agent to orchestrate tools in an interpretable sequence. The design follows a layered approach: the CLI provides the entry point, the Agno agent coordinates tool invocations, MCP servers expose analysis capabilities, and backend libraries handle low-level parsing and heuristics. This separation mirrors the principles of modern malware analysis platforms: multiple specialized stages, each with clear inputs and outputs, allowing for modular extension and failure isolation.

The architecture avoids the complexity of embedded graph embeddings (as in ~\cite{gemini2018}) or deep instrumentation (as required by full unpacking frameworks~\cite{polyunpack2007}), instead providing high-level abstractions that operationalize key insights from these approaches without implementing them end-to-end.

\begin{verbatim}
CLI -> Agno Agent -> MCP Servers -> Analysis Backends
                       |           (LIEF, Capstone, YARA, Androguard)
                       +-> OpenRouter LLM
\end{verbatim}

The agent receives high-level analysis goals and invokes tools via MCP. Each tool returns JSON-serializable results, which the agent summarizes for the user or passes to subsequent analysis stages. This design enables graceful degradation: if the LLM is rate-limited, local tools still function and produce structured reports.

\subsection{Tool Design Philosophy}
The tools are designed as high-level analysis primitives rather than raw CLI wrappers. Each tool returns structured, interpretable data that the agent can use to summarize risk. Examples:
\begin{itemize}[leftmargin=*]
  \item \texttt{extract\_crypto\_constants()} detects cryptographic constant sequences.
  \item \texttt{extract\_strings\_with\_context()} returns strings with byte context.
  \item \texttt{find\_suspicious\_syscalls()} flags syscall-related indicators.
  \item \texttt{extract\_permissions\_with\_risk()} assigns permission risk levels.
  \item \texttt{find\_hardcoded\_secrets()} detects API keys and tokens.
  \item \texttt{detect\_obfuscation\_techniques()} flags obfuscation heuristics.
\end{itemize}

\section{Implementation}
\subsection{Technology Stack}
\begin{tabular}{@{}ll@{}}
\toprule
Component & Technology \\ \midrule
Agent Framework & Agno \\ 
MCP Protocol & FastMCP \\ 
LLM Runtime & OpenRouter (free tier) \\ 
Binary Analysis & LIEF, Capstone, YARA (optional) \\ 
APK Analysis & Androguard (optional) \\ 
Language & Python 3.11 \\ 
Containerization & Docker, docker-compose, start.sh \\ \bottomrule
\end{tabular}

\subsection{Key Modules}
\begin{itemize}[leftmargin=*]
  \item \textbf{MCP servers}: static, dynamic, pattern, and APK services expose analysis tools.
  \item \textbf{Agent}: orchestrates tool calls and produces a concise risk summary.
  \item \textbf{CLI}: runs analysis in local mode or agent-orchestrated mode.
  \item \textbf{Reporting}: JSON report generation with timestamps.
\end{itemize}

\subsection{CLI Workflow}
The CLI exposes two primary usage modes. In local mode, the CLI runs all tools directly and produces a JSON report. In agent mode, the LLM summarizes the local results. In full orchestration mode, the agent can call tools directly and generate the final summary. This maps to three clear commands:
\begin{itemize}[leftmargin=*]
  \item Local analysis: \texttt{python -m agentic\_binary\_analysis.cli analyze --file <path> --kind binary|apk}
  \item Local analysis + LLM summary: add \texttt{--agent}
  \item Agent tool orchestration: add \texttt{--agent --agent-orchestrate}
\end{itemize}
This separation made it easy to test each step independently and keep the system usable even when the free-tier LLM rate-limits.

\subsection{Analysis Tools}
\textbf{Binary static analysis tools} form the foundation of the workflow. Rather than implementing feature extraction for graph-based similarity detection (as described in recent neural approaches~\cite{gemini2018}), these tools expose simpler but interpretable primitives: cryptographic constant sequences (RC4, AES SBOX patterns), string pools with surrounding byte context, and import/export enumeration. The crypto detection uses known key schedules and S-box patterns, allowing quick identification of cryptographic operations without requiring disassembly.

String extraction with context is deliberately designed to give analysts meaningful signals about functionality. A string like \texttt{accept} appearing near socket-related bytes suggests network functionality without requiring full semantic analysis. This approach trades depth for interpretability, keeping analysis understandable to both human analysts and LLM agents.

\textbf{Binary dynamic heuristics} focus on detecting anti-analysis and packing indicators. Rather than full instrumentation (as required by comprehensive unpacking frameworks~\cite{polyunpack2007}), these tools look for behavioral signatures: references to debugging APIs, checks for common analysis tools, and evidence of self-modifying code through entropy anomalies. Anti-analysis detection is critical because sophisticated malware often detects sandboxes and disables functionality when analyzed~\cite{polyunpack2007}. Syscall-related imports (\texttt{mmap}, \texttt{shmget}, \texttt{ptrace}) are flagged as potential process manipulation indicators.

\textbf{Pattern detection tools} estimate packing and obfuscation state through entropy analysis and YARA matching. Packed binaries typically exhibit high entropy ($> 6.0$ bits/byte) in data sections before unpacking occurs. Control-flow indicators—jump tables, indirect branches, and frequent relocation references—suggest dynamic code generation or polymorphism. When YARA is available, rule matching provides additional signals for known packing signatures.

\textbf{APK tools} extract permissions and assign risk levels based on security-sensitive groups (network, telephony, location, file access). Hardcoded secrets detection scans for API keys (Google Cloud, AWS patterns), cryptographic material, and authentication tokens. Network behavior analysis enumerates URLs embedded in resources and bytecode, which often reveals command-and-control infrastructure or exfiltration endpoints. Obfuscation detection flags class name scrambling, string encryption via constant pools, and DEX manipulation patterns—all of which suggest intentional disguise of functionality.

\textbf{Data formats and serialization} are carefully designed for agent interpretation. All tools return JSON with: file path, detected items (with confidence/context), and warnings. This structure allows the agent to post-process results, compare tools, and synthesize a risk narrative. The JSON format also enables reproducible logging and off-line analysis.

\subsection{Data Formats}
All tools return JSON-serializable dictionaries with consistent fields (file path, hits/matches, and any warnings). This makes it easy for the agent to summarize results and for the CLI to generate reports. The report format includes a UTC timestamp and the full analysis payload.

\subsection{Agent Orchestration}
The Agno agent receives the analysis objective and calls the MCP tools. It then summarizes the results in concise Markdown. For low-resource environments, the system supports a fallback path where tools run locally and the LLM only summarizes the JSON payload.

\subsection{Containerization}
The project includes a Dockerfile, docker-compose, and a start script. The start script validates required environment variables and warns if optional tools are missing. This keeps the workflow reproducible without forcing heavy dependencies on the host machine.

\subsection{Comparison with Related Approaches}
Recent work in binary analysis has explored neural network approaches for similarity detection~\cite{gemini2018} and comprehensive unpacking frameworks for hidden-code extraction~\cite{polyunpack2007}. The Gemini approach learns embeddings of binary functions to detect code reuse and malware variants across architectures. While powerful, such approaches require substantial training data, careful feature engineering, and computational resources for model inference.

PolyUnpack provides automated extraction and analysis of hidden code in self-modifying and runtime-packed binaries through bytecode instrumentation and memory capture. This comprehensive approach is highly effective for revealing obfuscated functionality but requires binary rewriting and system-level instrumentation.

In contrast, this implementation focuses on composable, lightweight heuristics that run locally on low-resource hardware without requiring model training or intrusive instrumentation. The tools are designed to be interpretable to both human analysts and LLM agents, enabling collaborative analysis. The trade-off is that this approach detects signals of obfuscation and suspicious behavior rather than completely extracting hidden code, but the interpretability and composability support agent-based reasoning and rapid feedback.

\subsection{Low-Disk Setup}
To support low-end hardware and minimal deployment footprint, a minimal dependency option was designed. The minimal requirements file (\texttt{requirements-min.txt}) installs only the agent framework, OpenRouter client, HTTP utilities, and test runner. Advanced analysis libraries (LIEF, Capstone, YARA, Androguard) remain optional and download-on-demand. This layered dependency approach supports multiple deployment scenarios: (1) constraint-heavy environments where disk is scarce, (2) rapid prototyping where not all analysis tools are needed, and (3) production environments where selective library loading is preferred. Even in minimal mode, the system still produces structured JSON reports and supports full agent orchestration through the LLM.

\section{Evaluation}
\subsection{Evaluation Goals and Design}
The evaluation aimed to validate three key claims: (1) the MCP tool abstraction remains effective for agent reasoning without requiring raw command wrappers, (2) the system can operate under resource constraints (minimal disk, rate-limited LLM), and (3) the workflow is reproducible via containerization. Rather than curating a large malware corpus (which would be infeasible in a low-resource environment), the evaluation uses synthetic samples to demonstrate tool correctness and workflow robustness.

\subsection{Test Protocol}
I validated each stage step-by-step: unit tests for core tool outputs, a CLI run without the agent, and an agent-orchestrated run. This layered approach isolates failures and allows incremental validation. Synthetic samples were chosen to: (a) avoid handling real malware on low-end hardware, (b) provide controlled test cases where expected outputs are known, and (c) demonstrate end-to-end tool chain functionality.

The evaluation protocol follows this sequence:
\begin{enumerate}
  \item \textbf{Unit tests}: Verify parsing and heuristic functions on minimal inputs.
  \item \textbf{Local analysis}: Run the CLI in local mode (no LLM) and validate JSON output structure.
  \item \textbf{Agent summary}: Run the CLI with \texttt{--agent} flag to verify LLM integration (using OpenRouter free tier).
  \item \textbf{Orchestration (optional)}: If rate limits allow, verify \texttt{--agent-orchestrate} mode where the agent directly calls MCP tools.
\end{enumerate}

\subsection{Test Validation}
\subsubsection{Core Tool Correctness}
Each analysis tool was validated on synthetic inputs with known properties:
\begin{itemize}[leftmargin=*]
  \item \textbf{String extraction}: Confirmed that \texttt{extract\_strings\_with\_context()} correctly identifies ASCII sequences and surrounding byte context.
  \item \textbf{Cryptographic constant detection}: Validated detection of RC4 S-boxes, AES key schedules, and common crypto patterns in binary stub code.
  \item \textbf{Entropy calculation}: Verified Shannon entropy computation on known data (e.g., repeated bytes yield low entropy $\approx 0$, uniform random data $\approx 8$ bits/byte).
  \item \textbf{Import enumeration}: Confirmed parsing of ELF and PE import tables (using LIEF when available).
\end{itemize}

\subsubsection{Binary Sample (\texttt{samples/sample.bin})}
The test binary is a minimal synthetic executable with known properties:
\begin{itemize}[leftmargin=*]
  \item Strings tool extracted \texttt{"hello\_world"} and surrounding context bytes, successfully demonstrating string pool parsing.
  \item Crypto constant detection found no cryptographic patterns (as expected for a benign stub).
  \item Suspicious syscall heuristics detected no anti-analysis indicators.
  \item Entropy calculation returned $\approx 3.16$ bits/byte, correctly indicating low compressibility and absence of encryption/packing.
  \item All tools returned valid JSON with no serialization errors.
\end{itemize}
\textbf{Result:} Low-risk classification confirmed; tool outputs validated as correct.

\subsubsection{APK Sample (\texttt{samples/sample.apk})}
The test APK is a synthetic Android package with planted indicators:
\begin{itemize}[leftmargin=*]
  \item Manifest parser correctly extracted declared permissions without errors.
  \item Hardcoded secrets scanner detected a Google API key pattern (\texttt{AIzaXXXXXXXXXXXXXXXXX}) in resource strings.
  \item URL detector found network endpoints (\texttt{https://example.com/api}, \texttt{https://analytics.example.org}).
  \item Obfuscation detection correctly identified no obfuscation in this synthetic sample (no class name scrambling, no bytecode encryption).
  \item Permissions risk assignment scored network permissions as ``high-risk'' and location as ``medium-risk'', aligning with security guidelines.
\end{itemize}
\textbf{Result:} High-confidence secret detection and consistent risk scoring; appropriate for alerting analysts to potential information exposure.

\subsubsection{Agent Reasoning Validation}
When running with the \texttt{--agent} flag, the Agno agent received structured JSON from local tools and produced human-readable summaries:
\begin{itemize}[leftmargin=*]
  \item The agent correctly identified the binary as benign based on low entropy and absent crypto patterns.
  \item For the APK, the agent flagged the hardcoded API key as a critical finding and highlighted network permissions as requiring further investigation.
  \item Agent output demonstrated interpretable reasoning (e.g., ``High entropy section suggests packing'' when entropy was elevated).
\end{itemize}

\subsection{Resource Constraint Testing}
\subsubsection{OpenRouter Free Tier Rate Limiting}
During full orchestration mode (\texttt{--agent-orchestrate}), the system encountered free-tier rate limits after 2-3 invocations. The fallback mechanism (local analysis + agent summary) functioned correctly and produced usable results within rate limits. This validates that the system remains operable even when LLM access is constrained.

\subsubsection{Low-Disk Mode}
Installation using \texttt{requirements-min.txt} succeeded on systems with limited disk space ($< 500$ MB free). The minimal installation skipped LIEF, Capstone, YARA, and Androguard. Tools that depend on these libraries gracefully degraded:
\begin{itemize}[leftmargin=*]
  \item Import enumeration returned empty results (no LIEF) but did not crash.
  \item String extraction and entropy analysis functioned normally using built-in Python libraries.
  \item YARA matching returned ``not available'' warnings but did not block analysis completion.
\end{itemize}
This confirms the modular dependencies design works as intended.

\section{Limitations}
This implementation operates within several deliberate constraints:
\begin{itemize}[leftmargin=*]
  \item \textbf{Heuristic-based detection}: The system detects signals of malicious behavior (high entropy, suspicious syscalls, API key exposure) rather than implementing comprehensive analysis as in graph embedding approaches~\cite{gemini2018} or instrumentation-based unpacking~\cite{polyunpack2007}. This reduces computational cost but sacrifices some depth of analysis.
  \item \textbf{Optional library dependencies}: Some advanced analyses depend on LIEF, Capstone, YARA, or Androguard. Systems lacking these libraries degrade gracefully but lose functionality. This is a deliberate trade-off to support low-disk environments.
  \item \textbf{Synthetic evaluation samples}: The evaluation uses synthetic test cases to demonstrate correctness. Real-world malware evaluation would require curated corpora and specialized handling (air-gapped systems, isolated networks). This is deferred to future work due to resource constraints.
  \item \textbf{OpenRouter free-tier rate limiting}: The system is designed to tolerate rate limits via fallback mode, but full LLM orchestration may be blocked under sustained load. Paid API tiers or local LLM integration would remove this constraint.
  \item \textbf{No dynamic analysis}: The implementation focuses on static analysis heuristics. Dynamic behavioral analysis (in emulators or sandboxes) is not included, which limits detection of runtime polymorphism and anti-analysis evasion as studied in~\cite{polyunpack2007}.
  \item \textbf{Architecture coverage}: Binaries are analyzed in architecture-agnostic ways (strings, entropy, heuristics), but detailed disassembly and architecture-specific semantics are not extracted. This limits false negatives but may increase false positives on complex architectures.
\end{itemize}

\section{Future Work}
Several extensions would increase the system's capability and depth:
\begin{itemize}[leftmargin=*]
  \item \textbf{Graph-based binary similarity}: Implement a simplified version of cross-platform similarity detection inspired by neural graph embedding~\cite{gemini2018}, using lightweight feature extraction optimized for LLM reasoning.
  \item \textbf{Automated unpacking}: Integrate selective memory capture and bytecode reconstruction techniques from unpacking frameworks~\cite{polyunpack2007}, targeting common packers (UPX, ASPack, themida).
  \item \textbf{Dynamic analysis integration}: Add support for lightweight emulation (e.g., QEMU user mode) or in-process hooking to capture runtime behavior.
  \item \textbf{Malware corpus evaluation}: Evaluate the system on real malware samples (in isolated environments) to measure false positive and false negative rates.
  \item \textbf{Richer control-flow analysis for binaries}: Extract and analyze call graphs, control-flow patterns, and interprocedural data flow using optional LIEF + custom graph analysis.
  \item \textbf{Expanded APK analysis}: Deep bytecode inspection, intent-based behavior analysis, and dynamic method resolution to detect intentional evasion.
  \item \textbf{Local LLM fallback}: Integrate local LLM frameworks (e.g., Ollama, vLLM) for air-gapped or offline operation without dependency on OpenRouter.
\end{itemize}

\section{Conclusion}
This project delivers an agentic workflow that automates binary and APK analysis using high-level MCP tools and LLM orchestration. The implementation prioritizes interpretability, composability, and low-resource operation, trading comprehensive analysis depth for clarity and scalability. The design successfully demonstrates that lightweight heuristics, when coordinated by an LLM agent, can provide meaningful risk signals for malware analysis without requiring specialized domain knowledge or heavy infrastructure.

The system's architecture separates concerns (CLI, agent, tools, backend) enabling modular extension and fault isolation. The tool abstraction layer avoids raw command wrappers, instead exposing semantic primitives (find crypto constants, extract secrets, measure entropy) that both human analysts and LLM agents can reason about. From a student perspective, this project demonstrates how combining established malware analysis techniques with modern LLM orchestration can yield a practical system that remains accessible in resource-constrained environments while supporting principled extension toward more sophisticated analysis as resources become available.

\begin{thebibliography}{2}

\bibitem{gemini2018}
Xiaojun Xu, Chang Liu, Qian Feng, Heng Yin, Le Song, and Dawn Song.
``Neural Network-based Graph Embedding for Cross-Platform Binary Code Similarity Detection.''
In \textit{Proceedings of the Network and Distributed Systems Security Symposium (NDSS)}, 2018.

\bibitem{polyunpack2007}
Paul Royal, Mitch Halpin, David Dagon, Robert Edmonds, and Wenke Lee.
``PolyUnpack: Automating the Hidden-Code Extraction of Unpack-Executing Malware.''
In \textit{Proceedings of the 16th USENIX Security Symposium}, 2007.

\end{thebibliography}

\end{document}
